\documentclass[Total.tex]{subfiles}

\begin{document}
	\chapter{Fields and Galois Theory}
		\section{Fields}
		
			\begin{Def}
				A \textbf{field} is a five-turple $(K,+,\times,0,1)$, where:
				
				\begin{itemize}
					\item $(K,+,0)$ is an additive group;
					
					\item $(K^\times,\times,0)$ is a multiplicative group; 
				\end{itemize}
			\end{Def}
			
			We have know the construction of
			
			\begin{gather*}
				Ring\rightarrow CRing\rightarrow ID\rightarrow Fields
			\end{gather*}
			
			\begin{proposition}
				Let $A$ be a domain. Then for every ring morphsim $\phi:A\rightarrow K$ where $K$ is a field, there exists factorization
				
				\begin{center}
					\begin{tikzcd}
						Quot(A)\arrow[r]&K\\
						A\arrow[u,hookrightarrow]\arrow[ru]
					\end{tikzcd}
				\end{center}
			\end{proposition}
			
			\begin{proof}
				It is well-defined, induced by
				
				\begin{gather*}
					\frac{a}{b}\mapsto \frac{\phi(a)}{\phi(b)}
				\end{gather*}
				
				Universal property omitted.	
			\end{proof}
			
			\subsubsection{Vector Spaces}
			
				\begin{lemma}
					Every $V\in Vect_k$ is free.
				\end{lemma}
				
				\begin{proof}
					Consider the Bases. (If infinite, rely on AC)
				\end{proof}
				
				\begin{corollary}
					Every exact sequence of modules over a field splits. 
				\end{corollary}
				
				\begin{proof}
					Each vector space is free, so is projective. And according to the proposition here: \ref{Homological-Algebra-Projective-Injective-Modules-Equivalent-Definition-of-Projective-Modules}
				\end{proof}
			
			\subsubsection{Characteristic}
			
				\begin{Def}
					Let $F$ be a field, and a ring map 
					
					\begin{gather*}
						\phi:\ZZ\rightarrow F\qquad 1\mapsto 1
					\end{gather*}
					
					Then, the kernel $ker(\phi)=(p)$, $char(F):=p$, is called the \textbf{characteristic} of $F$.
				\end{Def}
				
				Under this notation:
				
				\begin{Def}
					Let $F$ be a field:
					
					\begin{itemize}
						\item $char(F)=0,\QQ\subseteq F$
						\item $char(F)=p,F_p\subseteq F$
					\end{itemize}
					
					is called \textbf{prime field}, which is the smallest subfield of $F$.
				\end{Def}
			
			\subsection{Field Extensions}
			
				\begin{lemma}
					
					\label{Fields-Morphisms-from-Field-to-Ring}
					If $F$ is a field and $R$ is a nonzero ring, then, any ring homomorphism
					
					\begin{gather*}
						\phi:F\rightarrow R
					\end{gather*}
					
					is injective.
				\end{lemma}
				
				\begin{proof}
					Suppose $a\in Ker(\varphi);a\not=0$. Then, notice that
					
					\begin{gather*}
						\varphi(1)=\varphi(aa^{-1})=0\Rightarrow 1=\varphi(1)=0
					\end{gather*}
					
					Contradiction. 
				\end{proof}
				
				\begin{Def}
					If $F$ is a field contained in a field $E$, then $E$ is said to be a \textbf{field extension} of $F$. We shall wirte 
					
					\begin{gather*}
						E/F
					\end{gather*}
					
					to indicate that $E$ is an extension of $F$.
				\end{Def}
				
				
				\subsubsection{Category of Field Extension}
				
				With \ref{Fields-Morphisms-from-Field-to-Ring}, we have known: The ring morphsims between fields must be injective. And: Those morphisms represent the field extensions. 
				
				For morphism $\phi:F\rightarrow K$, $K$ can be seen as $F$-algebra.
				
				\begin{Def}
					Let $k$ be a field. There is a \textbf{category of field extensions} of $k$. 
					
					\begin{itemize}
						\item Objects: Morphisms $k\rightarrow E$;
						\item Morphsims: \begin{tikzcd}
							E\arrow[rr]&&E'\\
							&k\arrow[lu]\arrow[ru]
						\end{tikzcd}
						
						Take the notation of morphsims $Mor_k(E,E')$.
					\end{itemize}
				\end{Def}
				
				\begin{Def}
					(Tower of Extensions) A \textbf{tower of fields} $E_n/E_{n-1}/\dots/E_0$ consists of a sequence of extensions of fields $E_n/E_{n-1},\dots,E_1/E_0$.
				\end{Def}
				
				\begin{Def}
					
				\end{Def}
				
				\begin{lemma}
					内容...
				\end{lemma}
				
				\begin{proof}
					内容...
				\end{proof}
			
				\subsubsection{Finite Extensions}
				
				\begin{Def}
					Let $F/E$ be an extension of fields. The dimension of $F$ considered as an $E$-vector space is called \textbf{degree of extension} and is denoted as:
					
					\begin{gather*}
						[E:F]:=dim_E(F)
					\end{gather*}
				\end{Def}
				\subsubsection{Algebraic Extension}
				
					\begin{Def}
						Consider a field extension $F/E$. An element $\alpha\in F$ is said to be \textbf{algebraic} over $E$ is $\alpha$ is the root of some nonzero polynomial with coefficients in $E$.
						
						If all elements of $F$ are algebraic then $F$ is said to be an \textbf{algebraic extension} of $E$.
					\end{Def}
				
				\subsubsection{Minimal Polynomials}
				
				\subsubsection{Algebraic Closure}
			\subsection{Separation Extension}
	
	\chapter{The Fundamental Theorem of Galois Theory}
	
	\chapter{Algebraic Closure}
	
	\chapter{Infinite Galois Extensions}
	
\end{document}